\section{Introduction and scope}
\label{sec:introduction}
Counter-streaming electron populations provide a useful representation of the
strongly non-Maxwellian distributions generated during collisionless magnetic
reconnection. McGrae-Menge and collaborators~\cite{original} report that tagging
electrons by their initial side of a current sheet isolates sub-populations
whose intrinsic third moments are small compared with their combined heat
flux. Their exact moment decomposition explains how relative motion and
population asymmetry can generate a substantial aggregate third moment even
when each population has little intrinsic heat flux. Evolving additional
population moments is therefore physically motivated.

The present study examines two separate questions that this useful mechanism
does not settle: whether the numerical comparison starts from matched physical
states, and whether the implemented finite-relaxation closure has the kinetic
stability properties attributed to it. These are testable questions about
the identified source version and model, rather than an assertion that the
underlying mixture interpretation is wrong.

Our first result uses the authors' publicly archived initial moments and input
decks~\cite{dataset}. Identical nominal drift expressions enter a proper-momentum
initializer in PIC and an ordinary-velocity initializer in the fluid model.
The recorded initial moments show a current and pressure mismatch consistent
with that distinction. This is a reproducible control issue for interpreting
the nonlinear comparison; its eventual dynamical effect cannot be determined
from initial-state analysis alone.

Our second result is a controlled longitudinal stability test of the stated
continuum isotropization closure. A positive two-Maxwellian equilibrium is
stable to longitudinal electrostatic perturbations under the established
Penrose criterion~\cite{penrose1960}, but its
finite-relaxation fluid counterpart has an overstable mode. We give an explicit
linear system and check both its eigenvalue residual and Gauss constraint.
The test is homogeneous and nonrelativistic. It does not replace the source's
inhomogeneous reconnection calculation or show that the identified mode caused
any of its plotted features.

Two further results make the interpretation more precise. Continuity depends
on the divergence of the mass flux, so a zero normal velocity at a symmetry
plane does not prohibit density accumulation there. A conditioned third-moment
error bound also shows why a thermal-to-drift speed ratio alone does not give
a universal relative-error guarantee when population contributions cancel.
Neither result disputes the empirical parameter trend in the published scan.

\subsection{Relation to established work}
Heat-flux closure and its effect on collisionless response are established
subjects. Hammett and Perkins~\cite{hammett1990} developed fluid moment models
for Landau damping, while Wang and collaborators~\cite{wang2015} compared
multifluid moment models with PIC reconnection. Goldman and
collaborators~\cite{goldman2020} studied multibeam energy moments and the
distinction between component and aggregate centering. The mixture algebra
used below belongs to this established moment framework.

Closure-dependent instability spectra are also known. Ng and
collaborators~\cite{ng2019} demonstrate why a pressure-tensor model must be
checked against kinetic response rather than judged solely by its retained
variables. Our finite-relaxation example supplies a concrete test associated
with the closure and reference rate in the present source; we do not claim the
general phenomenon of closure-dependent unstable modes as new.

Work on heat-flux closure for nonlinear electrostatic two-stream systems
\cite{ingelsten2025,ingelsten2026} further limits a literal reading of broad
claims that no fluid closure can capture non-Maxwellian counter-streaming.
It does not establish prior use of the specific tagged-population closure in
electron-only reconnection. Our analysis therefore addresses explicit
initialization, continuity and stability statements, and preserves the
source's application-specific empirical evidence.

\subsection{Audited version and available evidence}
The supplied ten-page PDF is byte-identical to a fresh versioned download of
arXiv:2607.07106v1: 7,847,791 bytes, SHA-256
\texttt{fc5344e2ba8f4db98baf90d805bb6b7ce}\\
\texttt{8835c32ea7bf54bba4d62f405b955a3}.
The official history checked on 28 October 2026 lists the 8 July v1 and no
later revision or journal reference. The PDF displays 9 July, while the
experimental HTML displays a different compilation date. The original TeX
uses \verb|\date{\today}|; those compilation dates do not identify additional
scientific revisions. Appendix~\ref{app:source-evidence} records source and
data locators. No original author has been contacted, and no author
confirmation or private simulation-build access is claimed.
